\src{172}
such that $K'^{i}$ is a direct summand of $K^{i}$ for every $i$, gives rise to an
exact sequence of complexes of groups $\operatorname{Hom}(K'',L)$, and so on,
and hence to a connecting operator
\[
  H^{i}(K',L)\longrightarrow H^{i+1}(K'',L).
\]
The connecting operators associated with an exact sequence in $L$ are defined
similarly. The pairings (2.2) are compatible, in the usual sense, with these
connecting operators.

Now suppose that $\Ccat$ is a category in which every object $A$ admits an
injective resolution $C(A)$. One then finds, using one of the many variants of
the bicomplex theorem, that
\[
  H^{*}\bigl(C(A),C(B)\bigr)
  =H^{*}\bigl(\operatorname{Hom}(C(A),C(B))\bigr)
\]
is canonically isomorphic to
\[
  H^{*}\bigl(\operatorname{Hom}(A,C(B))\bigr)
  =\operatorname{Ext}^{*}(A,B).
\]
The connecting operators considered above give the connecting homomorphisms for
$\operatorname{Ext}$. Moreover, the pairings (2.2) here give associative
pairings compatible with the connecting homomorphisms:
\begin{equation}\tag{2.3}
  \operatorname{Ext}^{*}(A,B)\times\operatorname{Ext}^{*}(B,C)
  \longrightarrow\operatorname{Ext}^{*}(A,C).
\end{equation}
In particular, $\operatorname{Ext}^{*}(A,A)$ is an associative graded ring with
identity, and so on. (One proves similarly that the $\operatorname{Ext}$
functors operate on the derived functors of an arbitrary functor; we shall not
use this fact here.)

When the category in question is the category $\Ccat^{\Osh}$ of
$\Osh$-Modules on $X$, one therefore obtains pairings
\begin{equation}\tag{2.4}
  \operatorname{Ext}_{\Osh}^{p}(X;A,B)
  \times\operatorname{Ext}_{\Osh}^{q}(X;B,C)
  \longrightarrow\operatorname{Ext}_{\Osh}^{p+q}(X;A,C),
\end{equation}
which may be made explicit as indicated above. Moreover, the same developments,
with the category of abelian groups replaced by the category of sheaves of
abelian groups on $X$ and the functors $\operatorname{Hom}$ replaced by the
functors $\Homsh$, also define pairings having the same formal properties and,
this time, of a \fquote{local nature}:
\begin{equation}\tag{2.5}
  \Extsh_{\Osh}^{p}(A,B)\times\Extsh_{\Osh}^{q}(B,C)
  \longrightarrow\Extsh_{\Osh}^{p+q}(A,C).
\end{equation}
These are sharpened by observing that the homomorphisms (1.8) are compatible
with the pairings between the $\operatorname{Ext}$ groups.

